% Source of truth for the technical report. Compile with pdflatex.
\documentclass[11pt,a4paper]{report}
\usepackage[margin=25mm]{geometry}
\usepackage[T1]{fontenc}
\usepackage{booktabs}
\usepackage{hyperref}
\usepackage{longtable}
\usepackage{array}
\usepackage{amsmath}
\usepackage{xcolor}
\hypersetup{colorlinks=true,urlcolor=blue,linkcolor=blue}
\title{Compiler-Driven EPR Pre-Generation for Quantum Multi-Core Communication\\
\large Physical Validation, Memory Trade-offs, and Reproducible Evidence}
\author{BTP Research Project}
\date{September 2026}
\begin{document}
\maketitle
\begin{abstract}
This report studies whether a compiler that knows a future inter-core
communication trace can reduce quantum-communication service latency by
preparing request-specific EPR pairs before they are needed. The work starts
with topology-correct multi-core communication and analytical scheduling, then
executes the same compiler intent through two physical backends: ACE and native
SeQUeNCe. Every final numerical claim is tied to a hash-locked workload,
immutable contract, paired seeds, raw provenance, completion check, and
pair-identity audit.

This is a living technical report. Tables labelled ``audited'' may be cited;
tables labelled ``historical'' retain context but are not combined with newer
corrected studies. The six-lead fixed-policy matrices currently running are
described as pending and will be inserted only after their audits pass.
\end{abstract}
\tableofcontents
\clearpage

\chapter{Research Narrative and Contributions}

\section{Motivation}
Modular quantum computers require communication whenever a logical qubit moves
between cores. Teleportation replaces a direct quantum transfer with a shared
Bell pair, local operations, classical messages, and receiver correction. This
makes entanglement availability a scheduling resource. Pure on-demand
generation is robust but exposes generation latency after a request arrives.
Continuous-generation approaches maintain opportunistic stock, but do not know
which future request will consume a pair and can waste scarce coherence time.

The project asks a narrower, testable question: given a compiler trace, can a
runtime prepare the \emph{exact} pair for a future request early enough to help
without consuming the recovery capacity needed by demand traffic? The answer
cannot come from an offline schedule alone. It must include physical generation
failure, retries, memory lifetime, contention, pair identity, and fallback.

\section{Progression of the work}
\begin{longtable}{p{0.20\textwidth}p{0.69\textwidth}}
\toprule
Stage & Contribution and evidence boundary \\
\midrule
Topology and routing & Replaced an all-to-all classical-control assumption with
sparse mesh-aligned links, hop-by-hop forwarding, XY routing, and concurrent
multi-hop request handling. \\
Analytical policies & Built ODG, CGP, ACGP, fixed-lead, and dynamic-lookahead
schedulers with finite memory, generation capacity, storage-age fidelity, and
request-level traces. This stage explains scheduling opportunity but does not
model stochastic physical creation. \\
Native physical port & Connected compiler decisions to native SeQUeNCe
Barrett--Kok attempts, retry, decoherence, teleportation, swapping, and
receiver correction. \\
ACE physical integration & Connected request-specific preparation to ACE's
adaptive reservation/cache lifecycle and repaired compiler-pair provenance on
reused memory slots. \\
Controlled matrices & Locked a 4x4 QFT trace, conflict serialization, resource
profiles, paired seeds, raw output schema, and audits. The studies measure
readiness, service latency, fidelity at use, expiry, fallback, and rejection
pressure. \\
\bottomrule
\end{longtable}

\section{Relation to prior work}
Adaptive continuous entanglement generation motivates the CGP/ACGP baselines
\cite{Kolar2022}. AEPA motivates proactive allocation under limited quantum
memory \cite{AEPA}. SeQUeNCe provides the discrete-event quantum-network
simulation basis for the native implementation \cite{SeQUeNCe}. This project
does not claim to reproduce another paper's absolute numbers: the supplied QFT
trace, physical assumptions, control implementation, and latency endpoints are
explicitly frozen and reported here.

\section{Research paper study}
The literature review guided the experimental design rather than supplying
parameters by assumption. Four papers were studied alongside the SeQUeNCe
documentation and the project presentation prepared during the analytical
policy phase.

\begin{longtable}{p{0.23\textwidth}p{0.28\textwidth}p{0.38\textwidth}}
\toprule
Source & Main contribution & What this project adopted or tested \\
\midrule
Kolar et al. \cite{Kolar2022} & Continuous generation with an adaptive local
neighbour choice driven by request history. & Defined CGP and ACGP as
traffic-oblivious baselines. We retained finite cache capacity, expiry, and
history-based neighbour updates, then distinguished their opportunistic hits
from request-specific compiler readiness. \\
Zhan et al. \cite{Zhan2025} & Adaptive continuous generation implemented in
SeQUeNCe with protocol-level generation, decoherence, and purification. &
Motivated preserving physical protocol effects rather than treating a planned
generation as success. Purification remains outside the present compiler
comparison and is reported as a limitation. \\
Chen et al. \cite{AEPA} & Reinforcement-learning pre-allocation chooses
link-level inventory under traffic, TTL, occupancy, and expiry. & Motivated the
later question of how much capacity compiler work may occupy. The cap-four
wait-on-demand control isolates this capacity question before any learned
policy is attempted. \\
Suance et al. \cite{Suance2026} & ODG, CGP, and ACGP for multi-core meshes and
QFT, Cuccaro, Draper, and MCMT V-chain workloads. & Supplied the policy set,
mesh-oriented workload structure, and qualitative latency/fidelity trade-off
to reproduce analytically before physical compiler integration. \\
\bottomrule
\end{longtable}

\subsection{What the analytical reproduction established}
The policy-level reproduction implemented ODG, CGP, and ACGP on five mesh
sizes and four circuit-derived workloads. It reproduced the qualitative result
that pre-generation can lower exposed request latency while stored pairs reduce
fidelity. It did \emph{not} reproduce every published magnitude. In particular,
the reproduced fresh-pair ODG fidelity remained 1.0 while the paper's physical
baseline includes nonideal generation and control. The source paper does not
fully specify all gate timings, random-generation intervals, cache-discard
rules, purification triggers, or the authors' circuit/mapping configuration.
Accordingly, the analytical reproduction is evidence of policy behavior and
implementation coverage, not a numerical replication claim.

\subsection{How the paper study changed this project}
The review made three design decisions explicit. First, physical results must
retain stochastic generation, retry, memory noise, expiry, and teleportation
instead of translating a scheduler hit directly into lower latency. Second,
CGP/ACGP must be evaluated as speculative baselines with waste, rather than as
compiler-aware policies. Third, capacity and demand recovery need their own
controls: the fixed-lead and cap-four studies do not tune a policy merely to
outperform a baseline, but measure readiness, fidelity, expiry, and blocking
under identical resources.

\chapter{Problem and System Model}

\section{Research question}
Distributed quantum programs move logical qubits across cores. Teleportation
needs an EPR pair between participating cores. On-demand generation starts
pair generation after a request arrives. Compiler-driven pre-generation uses
known future requests to prepare exact request-specific pairs earlier. The
research question is whether that knowledge reduces service latency under
finite quantum memory, stochastic generation, retries, decoherence, expiry,
and contention.

The experiments replay communication requests rather than the complete
evolving 96-qubit QFT state. They execute physical generation, storage noise,
expiry, teleportation, and correction for every request. Raw ACE and native
SeQUeNCe milliseconds are not comparable because their latency endpoints and
physical paths differ. Policy comparisons are only valid within one backend.

\section{Controlled workload}
\begin{center}
\begin{tabular}{ll}
\toprule
Property & Value \\
\midrule
Topology & 4x4 mesh, 16 cores, 24 neighbour links \\
Logical qubits & 96 \\
Source layers & 766 \\
Transfers & 4,954 \\
Conflict-free sublayers & 2,831 \\
Trace SHA-256 & \texttt{61d97492195aea40fad48d5bc4e1b48b2dd019eea20fe24aada8c954e3073da1} \\
\bottomrule
\end{tabular}
\end{center}
Transfers sharing a core are serialized into conflict-free sublayers. A barrier
releases the next sublayer only after the prior sublayer terminates.

\section{Policies and memory model}
ODG generates after request arrival. Fixed compiler scheduling prepares the
exact pair two layers before use. Dynamic compiler scheduling selects the
latest planner-feasible preparation. CGP and ACGP speculatively generate
background pairs without compiler knowledge; ACGP adapts neighbour choice from
traffic history.

Historical static memory allocations were 3+1, 2+2, and 1+1 compiler/demand
partitions. The current shared pool has four physical slots per core with a
background or compiler occupancy cap of three. Demand can use any free slot;
endpoint allocation is atomic, and pending demand has priority over new
background work. Strict 4+0 disables fallback and measures compiler coverage.

\section{Definitions used throughout the report}
For a request $r$ released at trace layer $L_r$, fixed scheduling launches its
request-specific preparation at $L_r-\Delta$. A request is \emph{pre-ready}
only when its exact physical pair has completed before demand generation or
release begins. If no such pair exists, demand fallback is used. A compiler pair
used by another request is an identity failure, not a cache hit. Expiry/waste is
the fraction of compiler-created pairs that expire or remain unused at study
completion; it is a resource cost, not merely a failed request.

\chapter{Experimental Method and Reproducibility}

\section{Immutable experiment contract}
Each publishable cell freezes the trace hash, mesh, serialized layer map,
memory profile, scheduling parameters, physical runtime overrides, policy,
seeds, and metric definitions in JSON. The locked workload is a 4x4 mesh with
16 cores, 96 logical qubits, 766 source layers, 4,954 transfers, and 2,831
conflict-free serialized sublayers. The same integer seeds 0--29 are paired
within a backend/profile comparison.

\section{Audit gates}
Before a numerical table is promoted, raw \texttt{runs.json} must pass exact
trace validation, the complete seed set, request completion, generated/used/
expired pair accounting, unique compiler-pair use, target-identity checks, and
fidelity bounds. Summary CSV files and raw-artifact SHA-256 values are retained
as provenance. A seed-zero run is only a smoke test.

\section{Interpretation boundaries}
ACE and native SeQUeNCe answer complementary questions. ACE uses its adaptive
reservation/cache lifecycle and reports its application-service endpoint.
Native SeQUeNCe executes Barrett--Kok generation, retries, swapping,
teleportation, and receiver correction with a different endpoint. Raw
millisecond values are therefore never ranked across backends. The valid unit
of comparison is a paired policy direction and trade-off within the same
backend and memory profile.

\chapter{Engineering Contributions}

\section{Classical topology alignment}
The project replaced an all-to-all classical-router assumption with sparse
classical links that mirror the quantum topology. It added hop-by-hop
forwarding, XY mesh routing, serialized classical-channel timing, and request
identity and priority handling for concurrent multi-hop teleportation.
Linear, star, and 3x3 mesh scenarios passed 32 targeted regression tests.
The primary implementation commit is
\texttt{2cd7dd6bb5602b4caf052922b0285efda82087e8}.

\section{Compiler scheduling and native integration}
The trace parser validates topology, tracks logical-qubit placement, and builds
fixed-delta and dynamic-lookahead schedules. The initial analytical model
scheduled 2,358 fixed requests (47.60\%) and 2,408 dynamic requests (48.61\%).
Those values are scheduling evidence only because a scheduled generation was
treated as successful.

Native integration executes Barrett--Kok generation, detector failure, retry,
memory noise, expiry, teleportation, and receiver correction. An exact compiler
pair cannot be consumed by another request. An in-flight exact attempt is
adopted at request arrival; an absent exact pair uses demand fallback.

\section{Analytical policy framework}
The analytical stage introduced the policy vocabulary used in all later
experiments. ODG creates pairs only after demand. CGP periodically chooses a
neighbour uniformly when capacity is available. ACGP updates a local neighbour
probability distribution from recently used edges and penalizes neighbours that
hold unused pairs. Fixed compiler scheduling binds a preparation to one future
request at a constant lead; dynamic scheduling searches backward over a bounded
lookahead and chooses the latest analytically feasible layer.

This stage generated QFT, Cuccaro, Draper, and MCMT V-chain communication
workloads on mesh, torus, line, and tree topologies. It established qualitative
questions about latency, cache hits, storage-age fidelity, and waste. It is not
presented as a physical benchmark because it treats generation rounds as
successful and uses an analytical storage-fidelity model. Its role in this
report is to motivate the physical experiments and define falsifiable policy
expectations.

\section{Native physical request lifecycle}
For each compiler preparation, the native simulator reserves endpoint memory,
starts Barrett--Kok generation, and retries physical failures while the target
request remains unreleased. At request release it first consumes an exact ready
pair. In historical native studies an exact in-flight job may be adopted;
otherwise the request uses demand fallback. The resulting pair record includes
planned and actual times, target request identity, attempts, readiness, expiry,
and density-matrix fidelity. This lifecycle shows why an offline scheduled pair
is not automatically a physical latency reduction: failures, retries, and
queueing remain on the critical path.

\section{ACE lifecycle repair}
ACE could consume a compiler pair while retaining its compiler reservation,
timecard, and quota until nominal expiry. The stale state produced artificial
congestion. The repair releases compiler bookkeeping when the pair is used.
All final ACE static values below use the corrected lifecycle.

\chapter{Audited Results and Interpretation}

\section{Evidence status}
Results are separated by their evidence status. The shared-pool ACE fixed-lead
matrix is a completed corrected 30-seed study. The static and native
fixed-lead matrices are executing under newly locked contracts and will replace
the historical fixed-$\Delta=2$ rows below only after audit. Historical rows
remain useful for explaining the evolution of the system, but are not combined
with new values.

\section{Historical native physical 3+1 study}
\begin{center}
\begin{tabular}{lrrrr}
\toprule
Policy & Latency (ms) & Pre-ready & Fidelity & Expiry \\
\midrule
Full ODG, four demand memories & $0.2095 \pm 0.0045$ & -- & 0.9498 & -- \\
Matched ODG, one demand memory & $0.5214 \pm 0.0093$ & -- & 0.9498 & -- \\
Fixed compiler & $0.3135 \pm 0.0076$ & 42.18\% & 0.8630 & 0\% \\
Dynamic compiler & $0.3524 \pm 0.0083$ & 30.46\% & 0.9193 & 0\% \\
\bottomrule
\end{tabular}
\end{center}
Fixed gains more physical retry time. Dynamic launches later and uses fresher
pairs. This establishes a retry-time versus fidelity trade-off.

\section{Historical corrected ACE static matrix}
\begin{center}
\begin{tabular}{lrrr}
\toprule
Profile & Fixed latency / ready / fidelity & Dynamic latency / ready / fidelity & Dynamic vs fixed \\
\midrule
3+1 & 0.658419 / 51.87\% / 0.7682 & 0.575869 / 61.59\% / 0.8039 & 12.52\% faster \\
2+2 & 0.821357 / 32.97\% / 0.8496 & 0.741142 / 42.13\% / 0.8738 & 9.75\% faster \\
1+1 & 0.908687 / 22.85\% / 0.8267 & 0.910578 / 22.57\% / 0.7861 & no reliable winner \\
\bottomrule
\end{tabular}
\end{center}
The corrected static audit passed 300 cells, 1,486,200 request instances, and
407,712 compiler-pair records with zero audit errors. These fixed rows use a
single historical lead and are superseded, for fixed-policy conclusions, by the
six-lead study described next.

\section{Completed ACE shared-pool fixed-lead sensitivity}
The completed shared-pool ACE study uses four physical slots per core, compiler
occupancy cap three, demand fallback, and the six fixed leads
$\Delta\in\{1,2,3,4,5,6\}$. It contains 30 paired seeds per lead and 148,620
request instances per cell. The raw artifacts, immutable contracts, identity
audits, and SHA-256 provenance are indexed in the evidence appendix.

\begin{center}
\begin{tabular}{rrrrrr}
\toprule
$\Delta$ & Latency (ms) & Pre-ready & Fidelity at use & Expiry & Fallback \\
\midrule
1 & 0.598283 & 59.17\% & 0.8584 & 55.32\% & 40.80\% \\
2 & 0.635186 & 54.72\% & 0.7856 & 57.18\% & 45.26\% \\
3 & 0.726038 & 44.15\% & 0.7732 & 58.68\% & 55.82\% \\
4 & 0.746164 & 41.91\% & 0.7344 & 59.64\% & 58.09\% \\
5 & 0.784498 & 37.25\% & 0.7035 & 61.44\% & 62.75\% \\
6 & 0.818597 & 33.56\% & 0.6530 & 62.66\% & 66.44\% \\
\bottomrule
\end{tabular}
\end{center}

For this legacy shared-cap-three execution, shorter lead is better on all six
reported operational measures: it lowers latency, increases request-specific
readiness and fidelity at use, and reduces fallback and expiry. The result is
not evidence that a one-layer lead is universally optimal. It is conditional on
the trace, physical parameters, scheduler, and legacy demand admission. A
new atomic-admission cap-three control is required before attributing a future
cap-three versus cap-four difference solely to compiler capacity.

\section{Six-lead matrices in progress}
The report will add six-lead tables for ACE static 3+1, 2+2, and 1+1; native
static 3+1, 2+2, and 1+1; and native shared pool. The ACE shared-pool table
above also becomes the fixed-compiler evidence for the ACE adaptive-baseline
comparison. ODG, dynamic, CGP, and ACGP are controls, not lead sweeps.

\chapter{Validation and Current Work}

Audits check completion, pair conservation, unique pair use, memory bounds,
event ordering, timing-stage sums, and fidelity bounds. ACE focused verification
reported 25 passed and one skipped. Native adaptive integration tests reported
15 passed and one skipped.

The native ODG/CGP/ACGP full-trace seed-0 smoke completed 14,862 requests with
zero failures and passed its raw-pair audit. The 30-seed native adaptive matrix
is running during this report's initial preparation. Its values must not enter
a final comparison table until every seed completes and the pair audit passes.

\section{Next research steps}
\begin{enumerate}
\item Complete and audit native ODG/CGP/ACGP shared-pool seeds 0--29.
\item Compare native adaptive policies with native fixed and dynamic policies.
\item Evaluate static-bank CGP/ACGP and a retry-aware dynamic scheduler.
\item Vary capacity, lookahead, coherence, workload, and topology.
\end{enumerate}

\chapter{Evidence Index and Reproduction Guide}\label{chap:evidence}

\section{Canonical repositories}
\begin{itemize}
\item ACE source and contracts: \url{https://github.com/sanjeevkrishnaa/Compiler-Driven-ACE}
\item Native SeQUeNCe source and contracts: \url{https://github.com/sanjeevkrishnaa/SeQUeNCe}
\item ACE shared-pool fixed-lead report:
\url{https://github.com/sanjeevkrishnaa/Compiler-Driven-ACE/blob/codex/strict-compiler-4plus0/results/ace_fixed_lead_sensitivity_30seed_v1/report.md}
\item ACE fixed-lead scope and contracts:
\url{https://github.com/sanjeevkrishnaa/Compiler-Driven-ACE/blob/codex/strict-compiler-4plus0/FIXED_LEAD_STUDY_SCOPE.md}
\end{itemize}

\section{Local source-of-truth files}
\begin{longtable}{p{0.38\textwidth}p{0.51\textwidth}}
\toprule
Artifact & Role \\
\midrule
\texttt{MASTER\_README.md} & User-maintained project narrative and audited
Week 4 summary. \\
\texttt{SESSION\_CONTEXT\_HANDOFF.md} & Canonical scope, historical results,
and experiment constraints. \\
Week 1--3 Notion exports and Week 2 presentation & Engineering narrative,
literature-review notes, source-paper comparison, and historical validation.
Claims from these sources are reconciled against the current contracts and
audits before appearing as final results. \\
\texttt{experiments/*.json} & Immutable ACE contracts. \\
\texttt{results/ace\_fixed\_lead\_sensitivity\_30seed\_v1/} & Tracked
shared-pool report, CSV, raw hashes, contract hashes, and provenance. \\
\texttt{audit\_compiler\_results.py} & ACE completion and pair-identity audit. \\
\texttt{SeQUeNCe/docs/native-compiler-4x4-update.md} & Native physical
integration and historical 3+1 results. \\
\texttt{SeQUeNCe/example/multicore\_entanglement/contracts/} & Native
contracts, including the fixed-lead matrix. \\
\bottomrule
\end{longtable}

\section{Reproducing an audited ACE fixed-lead cell}
\begin{verbatim}
.venv/bin/python run_compiler_pregeneration.py \
  --trace /absolute/path/qft_requests.txt \
  --contract experiments/qft_4x4_shared_pool_fixed_delta2_v1.json \
  --profile shared-pool-4-cap3-fixed-delta2 \
  --output output/reproduction/delta2

.venv/bin/python audit_compiler_results.py \
  output/reproduction/delta2/runs.json \
  --expected-trace-sha256 61d97492195aea40fad48d5bc4e1b48b2dd019eea20fe24aada8c954e3073da1 \
  --output output/reproduction/delta2/audit.json
\end{verbatim}

\begin{thebibliography}{9}
\bibitem{Kolar2022}
Kolar et al., ``Adaptive, Continuous Entanglement Generation for Quantum
Networks,'' IEEE INFOCOM Workshop, 2022.
\bibitem{Zhan2025}
Zhan et al., ``Design and Simulation of the Adaptive Continuous Entanglement
Generation Protocol,'' IEEE QCNC, 2025.
\bibitem{AEPA}
Chen et al., ``Adaptive Entanglement Pre-Allocation for Low-Latency Quantum
Repeater Networks,'' IEEE INFOCOM Workshop, 2026.
\bibitem{Suance2026}
Suance et al., ``Adaptive Entanglement Management in Quantum Multi-Core
Architectures,'' ISLVSI manuscript, 2026.
\bibitem{SeQUeNCe}
X. Wu et al., ``SeQUeNCe: A Customizable Discrete-Event Simulator of Quantum
Networks,'' \emph{Quantum Science and Technology}, 2021.
\end{thebibliography}
\end{document}
